%! Author = sriadityavedantam
%! Date = 3/25/26

\lesson{50}{Self-Study Six}{Day 50}

\begin{definition}[Connectedness]
	A metric space $(X,d)$ is connected if it is impossible to write
	\[
		X = U_1\cup U_2
	\]
	as a disjoint union of nonempty open sets with
	\[
		U_1\cap U_2 = \varnothing.
	\]
\end{definition}

\begin{example}[Examples of Connected Sets]
	\[
		\r,\quad \r^n,\quad [a,b].
	\]
\end{example}

\begin{theorem}{
	Let $(X,d_X)$ and $(Y,d_Y)$ be metric spaces and let $f:X\to Y$ be continuous.
	If $X$ is connected, then $f(X)$ is connected.
}
	Suppose, for the sake of contradiction, that $f(X)$ is disconnected.
	Then
	\[
		f(X) = U_1\cup U_2
	\]
	where $U_1,U_2$ are nonempty disjoint open sets in $Y$.
	Then
	\[
		X = f^{-1}(U_1)\cup f^{-1}(U_2).
	\]
	Since $f$ is continuous, $f^{-1}(U_1)$ and $f^{-1}(U_2)$ are open.
	They are disjoint and nonempty.
	This contradicts connectedness of $X$.
	Hence $f(X)$ is connected.
\end{theorem}

\begin{theorem}[Intermediate Value Theorem]{
	Let $X$ be connected and $f:X\to\r$ be continuous.
	If $a,b\in f(X)$ and $a < r < b$, then $r\in f(X)$.
}
	Suppose, for the sake of contradiction, that $r\notin f(X)$.
	Let
	\[
		A \coloneqq (-\infty,r)
		\qquad\text{and}\qquad
		B \coloneqq (r,\infty).
	\]
	Then
	\[
		f(X)\subseteq A\cup B.
	\]
	So
	\[
		X = f^{-1}(A)\cup f^{-1}(B).
	\]
	Both sets are open, disjoint, and nonempty.
	This contradicts connectedness of $X$.
	Hence $r\in f(X)$.
\end{theorem}

\begin{definition}[Directional Derivative]
	Given $U\subset\r^n$ open, $f:U\to\r^m$, $a\in U$, and $u\in\r^n$, define
	\[
		D_u f(a) \coloneqq \lim_{t\to0}\frac{f(a + tu) - f(a)}{t}.
	\]
\end{definition}

\begin{remark}
	In single-variable calculus,
	\[
		f\uptick(a) = \lim_{t\to 0}\frac{f(a + t) - f(a)}{t}.
	\]
	Define $\lambda:\r\to\r$ by
	\[
		\lambda(t) \coloneqq f\uptick(a)t.
	\]
	Then
	\begin{align*}
		\lim_{t\to0}\frac{f(a+t) - f(a) - \lambda(t)}{t}
		&= \lim_{t\to0}\left(\frac{f(a+t) - f(a)}{t} - f\uptick(a)\right)\\
		&= 0.
	\end{align*}
	So for small $t$,
	\[
		f(a+t) - f(a) \approx f\uptick(a)t.
	\]
\end{remark}

\begin{definition}[Differentiability]
	Given $U\subset\r^n$ open and $a\in U$, a function $f:U\to\r^m$ is differentiable at $a$ if there exists a linear map $B:\r^n\to\r^m$ such that
	\[
		\frac{f(a+h) - f(a) - Bh}{\norm{h}}\to 0
		\qquad\text{as } h\to 0.
	\]
	In this case, $Bh$ approximates $f(a+h)-f(a)$.
\end{definition}

\begin{theorem}{
	If $f$ is differentiable at $a$, then for every $u\in\r^n$, the directional derivative $D_u f(a)$ exists and
	\[
		D_u f(a) = Bu.
	\]
}
	Since $f$ is differentiable at $a$,
	\[
		\frac{f(a+tu) - f(a) - B(tu)}{\norm{tu}}\to 0
		\qquad\text{as } t\to 0.
	\]
	Then
	\begin{align*}
		\frac{f(a+tu) - f(a)}{t}
		&= \frac{B(tu)}{t} + \frac{f(a+tu) - f(a) - B(tu)}{t}\\
		&= Bu + \frac{f(a+tu) - f(a) - B(tu)}{t}.
	\end{align*}
	The second term goes to $0$ as $t\to 0$.
	Hence
	\[
		\frac{f(a+tu) - f(a)}{t}\to Bu.
	\]
\end{theorem}

\begin{remark}
	We use $Df(a)$ to denote the derivative.
\end{remark}

\section{Lebesgue's Criterion for Riemann Integrability}

\begin{observation}
	The following is Thomae's function.
	It is continuous at every irrational point and discontinuous at every rational point.
	Nevertheless, it is Riemann integrable on $[0,1]$ with
	\begin{gather*}
	    \int_0^1 t(x)\,dx = 0.\\
	    t(x) =
		\begin{cases}
			1, & x=0,\\
			1/n, & x=m/n\in\q\setminus\{0\}\text{ in lowest terms},\\
			0, & x\notin\q.
		\end{cases}\\
	\end{gather*}
\end{observation}

\begin{definition}[Measure Zero]
	A set $A$ has measure zero if for all $\eps > 0$, there exists a countable collection of open intervals $\{O_n\}$ such that
	\[
		A\subseteq\bigcup_{n=1}^\infty O_n
		\qquad\text{and}\qquad
		\sum_{n=1}^\infty\abs{O_n}\leq\eps.
	\]
\end{definition}

\begin{theorem}[Lebesgue's Theorem]{
	Let $f$ be bounded on $[a,b]$.
	Then $f$ is Riemann integrable if and only if the set of discontinuities has measure zero.
}
	Let $M>0$ such that $\abs{f(x)}\leq M$ for all $x\in[a,b]$.
	Define
	\begin{gather*}
	    D \coloneqq \{x\in[a,b] : \text{$f$ is not continuous at $x$}\}.\\
	    D^\alpha \coloneqq \{x\in[a,b] : \text{$f$ is not $\alpha$-continuous at $x$}\}.\\
	\end{gather*}

	\textbf{($\impliedby$).}
	Suppose $D$ has measure zero.
	Let $\eps > 0$ and set
	\[
		\alpha \coloneqq \frac{\eps}{2(b-a)}.
	\]
	\begin{problem}
		Show there exists a finite collection of disjoint open intervals $\{G_1,\ldots,G_N\}$ covering $D^\alpha$ such that
		\[
			\sum_{n=1}^N \abs{G_n} < \frac{\eps}{4M}.
		\]
	\end{problem}
	\begin{problem}
		Let $K \coloneqq [a,b]\setminus \bigcup_{n=1}^N G_n$.
		Show $f$ is uniformly $\alpha$-continuous on $K$.
	\end{problem}

	\begin{problem}
		Construct a partition $P_\eps$ such that
		\[
			U(f,P_\eps)-L(f,P_\eps)\leq \eps.
		\]
	\end{problem}

	\textbf{($\implies$).}
	Suppose $f$ is Riemann integrable.
	Given $\eps > 0$ and $\alpha > 0$, choose a partition $P_\eps$ such that
	\[
		U(f,P_\eps)-L(f,P_\eps) < \alpha\eps.
	\]

	\begin{problem}
		(a) Show $D^\alpha$ has measure zero.

		(b) Deduce $D$ has measure zero.
	\end{problem}
\end{theorem}